\documentclass[10pt,a4paper]{amsart}
\usepackage[margin=19mm]{geometry}
\usepackage{amsmath,amssymb,mathtools}
\usepackage[scr=rsfs,cal=euler]{mathalfa}
\usepackage{xfrac}
\usepackage[unicode,hidelinks,bookmarks=false]{hyperref}
\newcommand{\ie}{i.e.\ }
\newcommand{\cf}{cf.\ }
\newcommand{\resp}{respectively\ }
\newcommand{\Subsection}{\subsection}
\providecommand{\nobreakdash}{\nobreak\hbox{-}}
\setlength{\emergencystretch}{2em}
\multlinegap0pt
\usepackage[shortlabels]{enumitem}
\usepackage{amssymb}
\usepackage{xcolor}
\newtheorem{lemma}{Lemma}
\title{Erd\H{o}s Rado Sunflower Theorem for Shifted Families}
\author{Tapas Kumar Mishra}
\date{Source: arXiv:2606.02667v2 (CC BY 4.0)}
\begin{document}
\maketitle

\noindent\emph{Source and licence.} Erd\H{o}s Rado Sunflower Theorem for Shifted Families. Version arXiv:2606.02667v2 (CC BY 4.0), distributed by arXiv under the Creative Commons Attribution 4.0 International licence, http://creativecommons.org/licenses/by/4.0/, as recorded in the arXiv licence metadata for this exact version and shown on the abstract page https://arxiv.org/abs/2606.02667v2. Full licence notice: \texttt{licenses/arxiv-2606.02667-CC-BY-4.0.txt}.\par\smallskip

\noindent\textit{Editorial scope.} Counted unit: the article's lemma lemma (source0.tex lines 344--351). The article's complete original proof of it is delivered with it (source0.tex lines 352--391, one \texttt{proof} environment of 40 lines). No mathematics is written, repaired or re-derived: the statement and the proof are the article's own lines, in the article's own notation, and no retained line is substituted or re-flowed.  No further source-local context is needed: every cross-reference of every retained line resolves inside the two delivered spans.  Nothing was omitted from any retained span, so the package carries no omission record.  No retained line contains a citation, so the wrapper carries no bibliography and no bibliography entry.  No retained line contains a figure environment, an image-inclusion command or drawing code, so no image is delivered and the package declares no asset dependency.  Editorial formatting only, in addition: an AMS \texttt{amsart} preamble that restates the article's own declarations and macros used by the retained lines, namely the environments \texttt{lemma} on separate counters, together with the standard packages those lines need.  The retained environments print in this document as Lemma 1 in order of appearance, whereas the article prints the same items with the numbering of its own full document.  The complete native source of the article as distributed in the arXiv e-print for this exact version is bundled unchanged as \texttt{sources/201921/source0.tex}.
\begin{lemma}
	Let $\mathcal{F}$ be a family of subsets of $[n]$ and let $1 \leq i, j \leq n$. Then the following hold:
	\begin{enumerate}[label=(\roman*)]
		\item For any $F \in \mathcal{F}$, $|\mathcal{C}_{ij}(F)| = |F|$.
		\item $|\mathcal{C}_{ij}(\mathcal{F})| = |\mathcal{F}|$.
		\item The matching number does not increase, i.e., $\nu(\mathcal{C}_{ij}(\mathcal{F})) \leq \nu(\mathcal{F})$.
	\end{enumerate}
\end{lemma}

\begin{proof}
	(i) The operator $\mathcal{C}_{ij}$ either leaves a set $F$ unchanged or replaces one element, $j$, with another, $i$. In both cases, the cardinality of the set is preserved.
	
	(ii) The map $F \mapsto \mathcal{C}_{ij}(F)$ is a bijection from $\mathcal{F}$ to $\mathcal{C}_{ij}(\mathcal{F})$. To see this, we only need to show it is injective. Suppose $\mathcal{C}_{ij}(F_1) = \mathcal{C}_{ij}(F_2) = G$. If neither set was shifted, $F_1 = F_2 = G$. If both were shifted, then $F_1 = (G \setminus \{i\}) \cup \{j\}$ and $F_2 = (G \setminus \{i\}) \cup \{j\}$, which implies $F_1 = F_2$. A case where one is shifted and the other is not (e.g., $G = (F_1 \setminus \{j\}) \cup \{i\}$ and $G=F_2$) leads to a contradiction with the definition of the shift, as the existence of $F_2 \in \mathcal{F}$ would have prevented $F_1$ from being shifted to $G$. Thus, the map is a bijection and $|\mathcal{C}_{ij}(\mathcal{F})| = |\mathcal{F}|$.
	
	Let $\mathcal{F}' = \mathcal{C}_{ij}(\mathcal{F})$. To show that $\nu(\mathcal{F}') \leq \nu(\mathcal{F})$, we will take a maximum matching in $\mathcal{F}'$ and construct a matching of the same size in $\mathcal{F}$.
	Let $\mathcal{M}' = \{B_1, B_2, \dots, B_p\}$ be a maximum matching in $\mathcal{F}'$, so $p = \nu(\mathcal{F}')$. The sets $B_k$ are pairwise disjoint. For each $B_k \in \mathcal{M}'$, there is a unique preimage $A_k \in \mathcal{F}$ such that $B_k = \mathcal{C}_{ij}(A_k)$.	
	A set $B_k$ is altered by the shift (i.e., $B_k \neq A_k$) only if $B_k = (A_k \setminus \{j\}) \cup \{i\}$. This implies that $i \in B_k$. Since the sets in the matching $\mathcal{M}'$ are pairwise disjoint, at most one set in $\mathcal{M}'$ can contain the element $i$. Therefore, \textbf{at most one set in $\mathcal{M}'$ could have been altered} by the $\mathcal{C}_{ij}$ operation.	
	We consider two cases:
	
	\textbf{Case 1: No set in $\mathcal{M}'$ was altered.}
	In this case, $B_k = A_k$ for all $k=1, \dots, p$. This means every set in $\mathcal{M}'$ is also in $\mathcal{F}$. Thus, $\mathcal{M}'$ itself is a matching of size $p$ in $\mathcal{F}$, so $\nu(\mathcal{F}) \ge p$.
	
	\textbf{Case 2: Exactly one set in $\mathcal{M}'$ was altered.}
	Without loss of generality, let $B_1$ be the altered set. Then:
	\begin{itemize}
		\item $B_1 = (A_1 \setminus \{j\}) \cup \{i\}$, where $A_1 \in \mathcal{F}$, $i \notin A_1$, $j \in A_1$, and $B_1 \notin \mathcal{F}$.
		\item For all $k \geq 2$, $B_k = A_k \in \mathcal{F}$, and since $B_1 \cap B_k = \emptyset$, we must have $i \notin B_k$.
	\end{itemize}
	
	Now, we construct a matching $\mathcal{M}$ in $\mathcal{F}$. Consider the collection $\{A_1, B_2,\allowbreak \dots, B_p\}$.
	If this collection is a matching, we are done. This occurs if $j \notin B_k$ for all $k \geq 2$, as $A_1$ is then disjoint from all $B_k$.
	Suppose the collection is not a matching. This can only happen if $A_1$ intersects some $B_k$ for $k \geq 2$. Since $A_1 = (B_1 \setminus \{i\}) \cup \{j\}$ and $B_1$ is disjoint from all other $B_k$, the intersection must be the element $j$. As the sets $\{B_k\}_{k \geq 2}$ are disjoint, at most one of them can contain $j$. Let this be $B_2$, so $j \in B_2$.
	Now we have $j \in B_2$ and $i \notin B_2$. Since $B_2$ was not altered by the shift (i.e., $\mathcal{C}_{ij}(B_2) = B_2$), the condition for shifting must have failed. This implies that the set $C_2 = (B_2 \setminus \{j\}) \cup \{i\}$ must already be in $\mathcal{F}$.
	
	Let us define a new collection of sets $\mathcal{M} = \{A_1, C_2, B_3, \dots, B_p\}$. We claim this is a matching of size $p$ in $\mathcal{F}$.
	\begin{enumerate}
		\item \textbf{Membership in $\mathcal{F}$}: $A_1 \in \mathcal{F}$ by definition. $C_2 \in \mathcal{F}$ as argued above. $B_k \in \mathcal{F}$ for $k \geq 3$ by assumption.
		\item \textbf{Pairwise Disjointness}:
		\begin{itemize}
			\item The sets $\{B_3, \dots, B_p\}$ are pairwise disjoint.
			\item For $k \geq 3$, $A_1 \cap B_k = \emptyset$ because $B_1 \cap B_k = \emptyset$ and $B_k$ contains neither $i$ nor $j$.
			\item For $k \geq 3$, $C_2 \cap B_k = \emptyset$ because $B_2 \cap B_k = \emptyset$ and $B_k$ contains neither $i$ nor $j$.
			\item Crucially, $A_1 \cap C_2 = \emptyset$. We know $A_1 = (B_1 \setminus \{i\}) \cup \{j\}$ and $C_2 = (B_2 \setminus \{j\}) \cup \{i\}$. Since $B_1$ and $B_2$ are disjoint, any intersection between $A_1$ and $C_2$ must involve $i$ or $j$. However, $i \notin A_1$ and $j \notin C_2$. Therefore, they are disjoint.
		\end{itemize}
	\end{enumerate}
	We have successfully constructed a matching $\mathcal{M}$ of size $p$ in $\mathcal{F}$.
	
	In all cases, the existence of a matching of size $p$ in $\mathcal{F}'$ implies the existence of a matching of size at least $p$ in $\mathcal{F}$. Therefore, $\nu(\mathcal{F}) \ge \nu(\mathcal{F}')$. 
\end{proof}

\end{document}
